
          %%%%% ~~~~~~~~~~~~~~~~~~~~ %%%%%

\chapter{Simultaneous diagonalizability}
\setcounter{theorem}{0}
\setcounter{equation}{0}

%\cite{vanDerVaart1996}
%\cite{Kosorok2008}

%\renewcommand{\theenumi}{\alph{enumi}}
%\renewcommand{\labelenumi}{\textnormal{(\theenumi)}$\;\;$}
\renewcommand{\theenumi}{\roman{enumi}}
\renewcommand{\labelenumi}{\textnormal{(\theenumi)}$\;\;$}

          %%%%% ~~~~~~~~~~~~~~~~~~~~ %%%%%

\section{Commuting operators on a complex vector space admit at least one simultaneous eigenvector}

          %%%%% ~~~~~~~~~~~~~~~~~~~~ %%%%%

\begin{proposition}
\mbox{}
\vskip 0.1cm
\noindent
A mutually commuting family of linear operators on a finite-dimensional
complex vector space admits at least one simultaneous eigenvector.
More precisely, suppose:
\begin{itemize}
\item
    $V$\, is a finite-dimensional vector space over the field $\C$ of complex numbers,
\item
    $\mathcal{F} \subset \End(V)$\, is a family of mutually commuting linear operators defined on \,$V$,\, i.e.,
    \begin{equation*}
    A_{1}\,A_{2} \;\; = \;\; A_{2}\,A_{1}\,,
    \quad
    \textnormal{for each \,$A_{1}, A_{2} \,\in\,\mathcal{F}$}
    \end{equation*}
\end{itemize}
Then, the following statements are true:
\begin{enumerate}
\item
    There exists non-zero vector \,$v \,\in\, V\,\backslash\{\,0\,\}$\,
    and a $\C$-valued function $\lambda : \mathcal{F} \longrightarrow \C$
    such that
    \begin{equation*}
    A(v) \; = \; \lambda({A}) \cdot v\,,
    \quad
    \textnormal{for each \,$A \in \mathcal{F}$}
    \end{equation*}
\item
    Furthermore, if the family \,$\mathcal{F} \subset \End(V)$\,
    is a vector subspace of \,$\End(V)$,\,
    then the above function
    \,$\lambda : \mathcal{F} \longrightarrow \C$\,
    is a linear function.
\end{enumerate}
\end{proposition}
\proof
\begin{enumerate}
\item
	We proceed by induction on \,$\dim(V)$.\,

	\vskip 0.3cm
	\noindent
	\textbf{Base step:}\; The Proposition is valid for $\dim(V) = 1$.
	\vskip 0.1cm
	\noindent
	If \,$\dim(V) = 1$,\, then any non-zero vector, say \,$v \in V\,\backslash\,\{\,0\,\}$,\,
	forms a basis for \,$V$,\, and every linear operator \,$T \in \End(V)$\, can be given by
	\,$v \longmapsto c_{T} \cdot v$,\, for some \,$c_{T} \in \C$.\,
	In particular, the non-zero vector \,$v$\, is thus a simultaneous eigenvector of \,$\mathcal{F}$.\,
	This establishes the base step.

	\vskip 0.3cm
	\noindent
	\textbf{Induction hypothesis:}\; The Proposition is valid for every complex vector space of dimension $< n$.

	\vskip 0.3cm
	\noindent
	\textbf{Induction step:}\;
	\vskip 0.1cm
	\noindent
	Suppose \,$1 < \dim(V) = n$.\,
	Observe that exactly one of the following two scenarios is true:
	\begin{itemize}
	\item
		Every operator \,$A \in \mathcal{F} \subset \End(V)$\,
		is a scalar multiple of the identity map
		\,$\id_{V} \in \End(V)$.\,
		\vskip -0.05cm
		\noindent
		The present Proposition is (trivially) true in this scenario.
	\item
		There exists \,$A \in \mathcal{F}\subset \End(V)$\,
		which is not a scalar multiple of the identity map
		\,$\id_{V} \in \End(V)$.\,
		\vskip -0.05cm
	    In this scenario, the Fundamental Theorem Algebra implies that \,$A$\,
	    admits at least one eigenvector
	    \,$w \in V\,\backslash\,\{\,0\,\}$,\, with associated eigenvalue \,$\beta \in \C$,\,
	    i.e., \,$A(w) = \beta \cdot w$.\,
	    We denote the eigenspace of the eigenvalue \,$\beta \in \C$\, by:
	    \begin{equation*}
	    E_{A}(\beta)
	    \;\; := \;\;
	    	\left\{\,
	    	    w^{\prime} \in V
	    	    \,\left\vert\;
	    	    A(w^{\prime}) \overset{{\color{white}1}}{=} \beta \cdot w^{\prime}
	    	    \right.
	    	    \right\}
	    \end{equation*}
	    Since \,$A$\, is not a scalar multiple of \,$\id_{V}$,\,
	    we see that \,$1 \,\leq\, \dim(\,E_{A}(\beta)\,) \,<\, n$\,
	    (otherwise, \,$\dim(\,E_{A}(\beta)\,) = n$\, would imply \,$A = \beta \cdot \id_{V}$).
	    By the induction hypothesis, the following family admits a simultaneous eigenvector:
	    \begin{equation*}
	    \mathcal{F}\,\vert_{E_{A}(\beta)}
	    \;\; := \;\;
	    	\left\{\,
	    	    B\,\vert_{E_{A}(\beta)} \in \End\!\left(\,\overset{{\color{white}.}}{E_{A}(\beta)}\,\right)
	    	    \,\left\vert\;
	    	    B \overset{{\color{white}1}}{\in} \mathcal{F}
	    	    \right.
	    	    \right\}
	    \end{equation*}
	    In other words, by the induction hypothesis, there exists non-zero
	    \,$v \in E_{A}(\beta) \subset V$\,
	    and a $\C$-valued function
	    \,$\mu : \mathcal{F}\,\vert_{E_{A}(\beta)} \longrightarrow \C$\,
	    such that
	    \begin{equation*}
	    B\vert_{E_{A}(\beta)}\!\left(\,\overset{{\color{white}.}}{v}\right)
	    \; = \;
	        \mu\!\left(\,B\vert_{E_{A}(\beta)}\,\right) \cdot v\,,
	    \quad
	    \textnormal{for each \,$B \in \mathcal{F}$}
	    \end{equation*}
	    Now, note that
	    \,$B\vert_{E_{A}(\beta)}\!\left(\,\overset{{\color{white}.}}{v}\,\right) \,=\, B(v)$,\,
	    since \,$v \in E_{A}(\beta)$.\,
	    Next, we define the $\C$-valued function
	    \,$\lambda : \mathcal{F} \longrightarrow \C$\,
	    by
	    \,$\lambda(B) \,:=\, \mu\!\left(\,B\vert_{E_{A}(\beta)}\,\right)$,\,
	    for each \,$B \in \mathcal{F}$.\,
	    Now, observe that:
	    \begin{equation*}
	    B(v)
	    \; = \;
	        B\vert_{E_{A}(\beta)}\!\left(\,\overset{{\color{white}.}}{v}\right)
	    \; = \;
	        \mu\!\left(\,B\vert_{E_{A}(\beta)}\,\right) \cdot v
	    \; = \;
	        \lambda(B) \cdot v\,,
	    \quad
	    \textnormal{for each \,$B \in \mathcal{F}$}
	    \end{equation*}
	    Thus, the non-zero vector \,$v \in V\,\backslash\,\{\,0\,\}$\,
	    is a simultaneous eigenvector of the family \,$\mathcal{F} \subset \End(V)$,\,
	    and the present Proposition is thus valid in this scenario.
	\end{itemize}
	This completes the proof of the induction step, and
    the desired result now follows by the principle of mathematical induction.
\item
    Simply observe that,
    for each \,$c_{1}, c_{2} \in \C$\, and each \,$A_{1}, A_{2} \in \mathcal{F} \subset \End(V)$,\,
    we have:
    \begin{eqnarray*}
    \lambda\left(\,c_{1}A_{1} \overset{{\color{white}.}}{+} c_{2}A_{2}\,\right) \cdot\, v
    & = &
        \left(\,c_{1}A_{1} \overset{{\color{white}.}}{+} c_{2}A_{2}\,\right)\!(\,v\,)
    \;\; = \;\;
        c_{1}A_{1}(v) \,+\, c_{2}A_{2}(v)
    \;\; = \;\;
        c_{1}\lambda(A_{1}) \cdot v \,+\, c_{2}\lambda(A_{2}) \cdot v
    \\
    & = &
        \left(\,c_{1}\lambda(A_{1}) \overset{{\color{white}.}}{+} c_{2}\lambda(A_{2})\,\right) \cdot\, v
    \end{eqnarray*}
    Since \,$v \neq 0$,\, it follows that
    \,$
    \lambda\left(\,c_{1}A_{1} \overset{{\color{white}.}}{+} c_{2}A_{2}\,\right)
    \,=\,
        c_{1}\lambda(A_{1}) \overset{{\color{white}.}}{+} c_{2}\lambda(A_{2})
    $,\,
    which proves the linearity of the function
    \,$\lambda : \mathcal{F} \longrightarrow \C$,\, as required.
    \qed
\end{enumerate}

          %%%%% ~~~~~~~~~~~~~~~~~~~~ %%%%%

          %%%%% ~~~~~~~~~~~~~~~~~~~~ %%%%%

